\begin{python0}
from solutions import *; clear()
\end{python0}
\sectionthree{When do I use the \texttt{while}-Loop?}

Suppose you want to write a program that adds numbers for a user but you
do not know in advance how many numbers to add. Of course you need to
give the user an opportunity to stop the adding process! There are two
ways of doing this. Either you can ask the user each time if he/she
wants to enter a number to add or to stop, or ... you can use a
\EMPHASIZE{sentinel value}. This is a special value entered by the user to
indicate that the processing loop should stop.

Let's try the simplest possible example. This program
prompts the user until he/she enters a sentinel value.
I'm choosing -1 to be the sentinel value.


\EMPHASIZE{i dont know how to properly format this right now}
prompt user for data
while value of data is not sentinel value:
process data
prompt user for data
Note that there is code duplication:
prompt user for data
while value of data is not sentinel value:
process data
prompt user for data

\EMPHASIZE{Removing Duplication in the sentinel \texttt{while}-loop}

Look at the sentinel while-loop again:

\EMPHASIZE{i dont know how to properly format this right now}
prompt user for data
while value of data is not sentinel value:
process data
prompt user for data

Remember that there is code duplication. If you want to get rid of code
duplication, you can do this:

while (1): // don't forget that 1 is type cast to true
prompt user for data
if value is the sentinel value:
break
process data

Note that some people do not like \texttt{break}s because
it's harder to tell when you get out of a loop because
you can issue a \texttt{break} statement anywhere in the body of the
\texttt{while} loop. They prefer the exit condition in the header of the
\texttt{while} loop so that visually one can tell quickly when you get out
of the loop.

Using the above technique, the following program:

\begin{console}[commandchars=\~\@\$]
int i = 0, sum = 0;

~textbf@std::cout << "gimme a number: "; // INPUT$
~textbf@std::cin >> i;$
while (~textbf@i != -1$) ~textbf@// CONDITION TO REPEAT LOOP$
{     
      sum += i;
      
      ~textbf@std::cout << "gimme a number: "; // INPUT$
      ~textbf@std::cin >> i;$
}
std::cout << "sum: " << sum << std::endl;
\end{console}
becomes:
\begin{console}[commandchars=\~\@\$]
int i = 0, sum = 0;
while (1)
{     
      ~textbf@std::cout << "gimme a number: "; // INPUT$
      ~textbf@std::cin >> i;$
      ~textbf@if (i == -1) break; // CONDITION TO EXIT LOOP$
      
      sum += i;
}
std::cout << "sum: " << sum << std::endl;
\end{console}
Make sure you study this program carefully!!!

